\section{Experimental Results and Analysis}
\label{sec:experimental-results}

This section evaluates the effectiveness of the proposed Dictionary-based
Attention (DA) block on visual object recognition and citation-network
classification. The visual experiments compare DA-HL with the conventional HL
baseline under three hypergraph construction strategies, namely $k$-NN,
clustering, and representation-based construction. The citation experiments
compare DA-DHGNN with DHGNN and established graph and hypergraph baselines. All
values in this section are the paper-reference results.

\subsection{Visual Object Recognition}

Table~\ref{tab:visual-main-results} summarizes the recognition accuracy on
ET-YaleB, MNIST, and RSSCN7. DA-HL consistently outperforms HL for all nine
dataset--generator combinations. On ET-YaleB, DA-HL obtains accuracies of
$80.9\%$, $80.4\%$, and $81.1\%$ with the $k$-NN, clustering, and
representation-based hypergraphs, respectively. These results correspond to
improvements of 1.7, 3.0, and 0.8 percentage points over HL. The largest gain
on this dataset is obtained with clustering, whose relatively coarse
hyperedges are more likely to contain weakly related samples.

On MNIST, DA-HL reaches $69.5\%$, $67.7\%$, and $69.9\%$ under the three
construction strategies, improving the corresponding HL results by 2.3, 3.4,
and 1.5 percentage points. The 3.4-point improvement for clustering is the
largest gain among all visual experiments. On RSSCN7, the improvements are
smaller but remain consistent: DA-HL increases accuracy from $67.5\%$ to
$68.6\%$ for $k$-NN, from $66.4\%$ to $68.1\%$ for clustering, and from
$68.1\%$ to $68.7\%$ for representation-based construction.

\begin{table}[t]
\centering
\caption{Visual object recognition accuracy (\%). The change is measured in
percentage points relative to HL.}
\label{tab:visual-main-results}
\scriptsize
\begin{tabular}{llrrr}
\toprule
Dataset & Generator & HL & DA-HL & Change \\
\midrule
\multirow{3}{*}{ET-YaleB}
 & $k$-NN & 79.2 & 80.9 & +1.7 \\
 & Clustering & 77.4 & 80.4 & +3.0 \\
 & Representation & 80.3 & \textbf{81.1} & +0.8 \\
\midrule
\multirow{3}{*}{MNIST}
 & $k$-NN & 67.2 & 69.5 & +2.3 \\
 & Clustering & 64.3 & 67.7 & +3.4 \\
 & Representation & 68.4 & \textbf{69.9} & +1.5 \\
\midrule
\multirow{3}{*}{RSSCN7}
 & $k$-NN & 67.5 & 68.6 & +1.1 \\
 & Clustering & 66.4 & 68.1 & +1.7 \\
 & Representation & 68.1 & \textbf{68.7} & +0.6 \\
\bottomrule
\end{tabular}
\end{table}

The results reveal two consistent patterns. First, representation-based
construction produces the highest absolute accuracy on all three datasets,
indicating that its initial hypergraph is more reliable than those generated by
$k$-NN or clustering. Second, the relative benefit of DA is inversely related
to the quality of the initial hypergraph. Clustering obtains the largest gains
of 1.7--3.4 points, followed by $k$-NN with gains of 1.1--2.3 points, whereas
the representation-based method gains 0.6--1.5 points. This behavior supports
the interpretation of DA as a noise-suppression module: when a hyperedge
contains more anomalous or weakly related vertices, sample-specific membership
weights have a greater effect.

The average improvement on RSSCN7 is 1.1 points, compared with 1.8 points on
ET-YaleB and 2.4 points on MNIST. RSSCN7 uses 2,048-dimensional ResNet features,
while ET-YaleB and MNIST use vectorized pixels. The stronger semantic features
produce a cleaner initial neighborhood structure, leaving less noise for the
DA block to correct.

\subsection{Citation-Network Classification}

Table~\ref{tab:citation-main-results} compares DA-DHGNN with classical graph
and hypergraph models on Cora, Citeseer, and Pubmed. DA-DHGNN achieves the best
accuracy on all three datasets: $83.8\%$ on Cora, $72.1\%$ on Citeseer, and
$81.5\%$ on Pubmed. Relative to the DHGNN baseline, the improvements are 1.3,
0.9, and 1.7 percentage points, respectively. DA-DHGNN also exceeds the best
non-DHGNN competitor by 1.6 points on Cora, 0.8 points on Citeseer, and 1.2
points on Pubmed.

\begin{table}[t]
\centering
\caption{Citation-network classification accuracy (\%).}
\label{tab:citation-main-results}
\scriptsize
\begin{tabular}{lrrr}
\toprule
Method & Cora & Citeseer & Pubmed \\
\midrule
Planetoid & 75.7 & 64.7 & 77.2 \\
DeepWalk & 67.2 & 43.2 & 65.3 \\
ICA & 75.1 & 69.1 & 73.9 \\
MoNet & 81.7 & -- & 78.8 \\
Chebyshev & 81.2 & 69.8 & 74.4 \\
ManiReg & 59.5 & 60.1 & 70.7 \\
LP & 68.0 & 45.3 & 63.0 \\
GCN & 81.5 & 70.3 & 79.0 \\
HGNN & 81.6 & 70.7 & 80.1 \\
HLPN & 82.2 & 71.3 & 80.3 \\
DHGNN & 82.5 & 71.2 & 79.8 \\
DA-DHGNN & \textbf{83.8} & \textbf{72.1} & \textbf{81.5} \\
\bottomrule
\end{tabular}
\end{table}

The improvement on citation networks is smaller than the largest gain observed
in visual recognition. Citation relations provide an explicit structure that
is generally more reliable than a hypergraph inferred solely from visual
features. Consequently, fewer anomalous memberships need to be suppressed.
Nevertheless, the consistent gains show that reweighting vertices within each
hyperedge remains beneficial even when the input relations are explicitly
observed.

\subsection{Robustness and Parameter Analysis}

The robustness experiments further demonstrate that DA becomes more effective
as the initial hypergraph deteriorates. On ET-YaleB, increasing the $k$-NN
neighborhood size from 5 to 40 reduces HL accuracy from $79.2\%$ to $61.5\%$,
whereas DA-HL decreases from $80.9\%$ to $72.9\%$. The corresponding advantage
of DA-HL therefore expands from 1.7 to 11.4 percentage points. A similar trend
is observed after adding Gaussian noise. At variance 0.8, HL falls to $39.9\%$,
while DA-HL retains $66.1\%$ accuracy, yielding a 26.2-point improvement.

On Cora, progressively removing hyperedge connections increases the advantage
of DA-DHGNN. At removal rates of 0.0, 0.2, 0.4, 0.6, and 0.8, the gains over
DHGNN are 1.3, 2.1, 3.1, 6.9, and 6.9 percentage points, respectively. These
results confirm that the attention block improves robustness by limiting the
propagation of unreliable memberships.

The parameter analysis shows that performance is relatively insensitive to the
dictionary size $K$. Across $K\in\{10,50,100,200,500\}$, the reported Cora
accuracy remains between $81.3\%$ and $83.8\%$, Pubmed remains between $79.2\%$
and $81.5\%$, and Citeseer remains between $71.1\%$ and $72.1\%$. In contrast,
the sparsity parameter $\alpha$ has a stronger influence and should be tuned for
each dataset. Overall, these experiments show that DA provides its largest
benefit when the original hypergraph contains noisy, excessive, or missing
connections, while requiring only moderate sensitivity to dictionary size.
